\documentclass[12pt,twocolumn]{article}
\usepackage[margin=2cm,columnsep=0.5cm]{geometry}
\usepackage{amsmath,amssymb}
\usepackage{hyperref}
\usepackage{xcolor}

\definecolor{bctnavy}{RGB}{11,37,69}
\definecolor{bctblue}{RGB}{13,71,161}

\title{\textbf{\color{bctnavy}BCT Appendix AH1: The Lost Archive}\\
{\large\color{bctblue}Complete Dataset Specifications and Processing\\
Parameters for Seven BCT SAR Archaeological Sites}}
\author{Michel Robert Cabri\'{e}\\
{\small Independent Researcher, Victoria, Australia}\\
{\small ORCID: 0009-0007-9561-9859}}
\date{23 March 2026}

\begin{document}
\twocolumn[
\maketitle
\begin{center}\rule{0.9\textwidth}{0.4pt}\end{center}
\begin{center}
\parbox{0.88\textwidth}{\small\textbf{Abstract.}
This appendix provides complete technical specifications for
BCT Bessel ratio SAR processing of seven archaeological
mystery sites described in BCT Letter~116. For each site
we give: exact coordinates, recommended dataset parameters,
processing code (Google Earth Engine), expected anomaly
signatures, and cross-validation approach. All datasets
are publicly available at zero cost.
\smallskip\hrule\smallskip}
\end{center}
\vspace{4pt}
]

\section{AH1.1 --- BCT Bessel Ratio Filter: Technical Definition}

The BCT Bessel ratio filter maps SAR backscatter intensity
$I(x,y)$ to a dimensionless anomaly field:

\begin{equation}
\mathcal{B}(x,y) = I(x,y) \times \frac{j_{1,1}}{j_{2,1}}
= I(x,y) \times \frac{2.4048}{3.8317} = 0.6276 \cdot I(x,y)
\end{equation}

The anomaly significance metric is:
\begin{equation}
Z(x,y) = \frac{\mathcal{B}(x,y) - \mu_\mathcal{B}(r)}
{\sigma_\mathcal{B}(r)},
\end{equation}
where $\mu_\mathcal{B}(r)$ and $\sigma_\mathcal{B}(r)$ are
the mean and standard deviation of $\mathcal{B}$ computed
within a context radius $r$ (typically 50~km). Detection
threshold: $Z > 3$ (3-sigma significance).

The physical basis: structures encoding the BCT void ratio
$\mathcal{R} = r_\mathrm{tet}/r_\mathrm{oct} = 0.5426$
produce resonant SAR backscatter at the $j_{1,1}/j_{2,1}$
Bessel ratio. The filter is sensitive to this geometric
signature regardless of the surface material --- stone,
mud-brick, timber, or bedrock all produce the same
Bessel ratio signature if their void geometry matches
BCT predictions.

\section{AH1.2 --- Google Earth Engine Code Template}

The following GEE JavaScript template applies the BCT
Bessel ratio filter to any of the seven sites. Substitute
the coordinates and date range as specified in
Sections~AH1.3--AH1.9.

\begin{verbatim}
// BCT Bessel Ratio SAR Filter
// Substitute SITE_LON, SITE_LAT, START, END

var j11 = 2.4048;
var j21 = 3.8317;
var bct_ratio = j11 / j21; // 0.6276

var site = ee.Geometry.Point([SITE_LON, SITE_LAT]);
var roi = site.buffer(50000); // 50km context

var s1 = ee.ImageCollection('COPERNICUS/S1_GRD')
  .filterBounds(roi)
  .filterDate('START', 'END')
  .filter(ee.Filter.eq('instrumentMode','IW'))
  .select('VV');

var median = s1.median().clip(roi);

// Apply BCT Bessel ratio filter
var bct_b = median.multiply(bct_ratio);

// Compute anomaly (Z-score)
var mean = bct_b.reduceRegion({
  reducer: ee.Reducer.mean(),
  geometry: roi, scale: 10
}).getNumber('VV');
var std = bct_b.reduceRegion({
  reducer: ee.Reducer.stdDev(),
  geometry: roi, scale: 10
}).getNumber('VV');

var z_score = bct_b.subtract(mean).divide(std);
var anomaly = z_score.gt(3); // 3-sigma detection

Map.centerObject(site, 12);
Map.addLayer(z_score, {min:-2,max:5,
  palette:['blue','white','red']},
  'BCT Bessel Ratio Z-score');
Map.addLayer(anomaly.selfMask(),
  {palette:['red']}, 'BCT Anomaly >3sigma');
\end{verbatim}

\section{AH1.3 --- Site 1: Hawara Labyrinth}

\textbf{Coordinates:} 29.2667°N, 30.9000°E\\
\textbf{GEE Point:} \verb|[30.90, 29.27]|\\
\textbf{Buffer:} 10~km (focused on pyramid precinct)\\
\textbf{Date range:} 2017--2026 (maximum scenes)\\
\textbf{Primary band:} VV polarisation\\
\textbf{Expected signature:} Grid pattern of wall reflections
at 8--12~m depth, detectable through shallow soil via
differential subsurface penetration at C-band (5.6~cm).

\textbf{Cross-validation:} Compare BCT anomaly map with
published 2008 Mataha Expedition GPR grid showing walls
at 8--12~m depth south of pyramid. Spatial correspondence
at $<$100~m constitutes confirmation.

\textbf{Biondi coordination:} When Dr Biondi publishes
HarmonicSAR Hawara results (expected Summer 2026), apply
BCT Bessel ratio to his published radar echo strength maps
for independent validation.

\textbf{Special note on central object:} The 40~m anomalous
object at Hawara centre. If metallic (as some scans suggest),
it will produce an extremely strong BCT Bessel ratio anomaly
due to high radar reflectivity. Expected Z-score $>10$.

\section{AH1.4 --- Site 2: Richat Structure (Atlantis)}

\textbf{Coordinates:} 21.1167°N, 11.3667°W\\
\textbf{GEE Point:} \verb|[-11.37, 21.12]|\\
\textbf{Buffer:} 30~km (full ring structure)\\
\textbf{Date range:} 2015--2026\\
\textbf{Primary band:} VV + VH (dual polarisation)

\textbf{Discriminating natural vs plasma origin:}
Natural eroded dome: random backscatter, no Bessel ratio
structure. Plasma discharge scar (Lichtenberg origin):
concentric Bessel ratio rings matching $j_{1,1}$, $j_{2,1}$,
$j_{3,1}$ spacing ratios 1:1.594:2.295.

\textbf{BCT anomaly prediction:} If Richat is a plasma
discharge scar, the ring spacing ratio should satisfy:
\begin{equation}
\frac{R_2}{R_1} = \frac{j_{2,1}}{j_{1,1}} = 1.594,
\end{equation}
where $R_1$ and $R_2$ are the inner and middle ring radii.
Observed ring radii: inner $\approx$5~km, middle $\approx$10~km,
outer $\approx$20~km. Ratio inner/middle = 0.50, close to
$j_{1,1}/j_{2,1} = 0.628$ but not matching. This suggests
natural dome origin is more likely --- but BCT SAR will
settle it definitively.

\textbf{Flood evidence:} Apply BCT Bessel ratio to the
southwest sand-drift region to characterise flood direction
and magnitude.

\section{AH1.5 --- Site 3: Alexandria Harbour (Cleopatra)}

\textbf{Coordinates:} 31.2000°N, 29.9000°E\\
\textbf{Method:} GEBCO bathymetry (0.25 arcminute,
available at \url{https://www.gebco.net})\\
\textbf{Processing:} BCT Bessel ratio applied to depth
gradient $|\nabla d(x,y)|$ rather than SAR backscatter.

Structured architectural ruins (rectangular rooms, column
bases, canal walls) produce gradient anomalies at right
angles inconsistent with natural sedimentation. BCT Bessel
ratio filter amplifies the spatial frequency corresponding
to OHC-resonant room geometry.

\textbf{Known structures:} Marine archaeology has located
several large stone blocks and column bases in the eastern
harbour at 5--8~m depth. BCT Bessel ratio will reveal
whether these are isolated finds or part of a systematic
geometric grid (library / palace complex).

\section{AH1.6 --- Site 4: Persian Gulf (Garden of Eden)}

\textbf{Coordinates:} 27°N, 53°E (central Gulf)\\
\textbf{Method:} GEBCO bathymetry + Sentinel-1 coastal SAR\\
\textbf{Target depth:} 0--40~m (pre-flood river valley)

\textbf{The Ur-Shatt River:} Ancient river systems are
visible in Gulf bathymetry at 20--40~m depth. BCT Bessel
ratio applied to bathymetric gradient will map the
pre-flood landscape and identify any anomalous geometric
features (settlement mounds, canal networks) inconsistent
with natural river valley morphology.

\textbf{Flood timeline:} Sea level rose $\sim$60~m between
15,000 and 6,000~BCE. Final Gulf flooding: $\sim$6,000~BCE.
Any settlement predating 6,000~BCE would be at 0--40~m
depth, within GEBCO resolution.

\section{AH1.7 --- Site 5: Mount Ararat (Noah's Ark)}

\textbf{Coordinates:} 39.7020°N, 44.2980°E\\
\textbf{Elevation target:} 4,700--4,900~m (anomaly zone)\\
\textbf{Dataset:} Sentinel-1 IW mode, mountain geometry\\
\textbf{Date filter:} Late summer (minimum snow cover)

\textbf{Anomaly signature:} A large wooden structure under
permafrost produces a Bessel ratio anomaly at the organic-
inorganic interface. Timber decays into lignin-rich matrix
with characteristic dielectric properties distinct from
surrounding basalt. Expected anomaly shape: elongated
($\sim$170~m), oriented NW-SE (consistent with published
satellite photographs of the anomaly).

\textbf{Falsification test:} If no BCT anomaly present
in late-summer low-snow Sentinel-1 data at the published
anomaly coordinates, the Ararat object is either natural
basalt or too deeply buried for C-band penetration.

\section{AH1.8 --- Site 6: Antarctic Peninsula}

\textbf{Coordinates:} 67°S, 65°W (primary target zone)\\
\textbf{Dataset:} BEDMAP3 (British Antarctic Survey,
freely available)\\
\textbf{Processing:} BCT Bessel ratio applied to radar
echo strength in BEDMAP3 radargrams.

\textbf{Key distinction:} Natural bedrock topography
(river valleys, glacially carved fjords) has characteristic
fractal geometry with no preferred Bessel ratio. Any
structured surface feature from the pre-glaciation
period (before $\sim$12,000~BCE) would appear as a
non-fractal anomaly at BCT Bessel ratio significance.

\textbf{Priority zone:} The Antarctic peninsula between
65°S and 70°S, where bedrock lies within 1~km of the
surface and pre-glaciation surface features are most
likely to be preserved.

\section{AH1.9 --- Site 7: Nazca Plateau}

\textbf{Coordinates:} 14.7300°S, 75.1300°W\\
\textbf{Dataset:} Sentinel-1 + SRTM elevation\\
\textbf{Processing:} BCT Bessel ratio applied to surface
SAR backscatter AND to subsurface GPR-derived void maps.

\textbf{Angular geometry test:} The Nazca lines include
geometric figures (spirals, trapezoids, animal outlines)
whose angular relationships BCT predicts should encode
Bessel mode ratios. Specifically:
\begin{equation}
\frac{\theta_1}{\theta_2} \approx \frac{j_{1,1}}{j_{2,1}}
= 0.628 \text{ for paired line sets}.
\end{equation}
This is a non-invasive angular analysis of published
Nazca line maps requiring no new data.

\textbf{Puquio system:} The underground water channels
(puquios) beneath the plateau have never been fully mapped.
BCT Bessel ratio SAR will detect their geometry as linear
anomalies in subsurface-sensitive C-band backscatter.

\section{AH1.10 --- Summary Table}

\begin{center}
{\small
\begin{tabular}{p{1.8cm}p{2.2cm}p{2.0cm}p{1.5cm}}
\hline\hline
Site & Primary Dataset & BCT Signature & Cost \\
\hline
Hawara & Sentinel-1 SAR & Grid walls >3$\sigma$ & Zero \\
Richat & Sentinel-1 SAR & Ring spacing ratio & Zero \\
Alexandria & GEBCO bathy & Rect. gradient & Zero \\
Persian Gulf & GEBCO bathy & River anomaly & Zero \\
Ararat & Sentinel-1 SAR & Elongated void & Zero \\
Antarctica & BEDMAP3 & Non-fractal bedrock & Zero \\
Nazca & Sentinel-1 SAR & Angular Bessel & Zero \\
\hline
\textbf{Total} & \textbf{Open data} & \textbf{7 tests} & \textbf{Zero} \\
\hline\hline
\end{tabular}
}
\end{center}

\section*{Data Availability}

All datasets referenced in this appendix are freely available:
\begin{itemize}
\item Sentinel-1 SAR: \url{https://scihub.copernicus.eu}
\item Google Earth Engine: \url{https://earthengine.google.com}
\item GEBCO Bathymetry: \url{https://www.gebco.net}
\item BEDMAP3: \url{https://www.bas.ac.uk/project/bedmap}
\item SRTM Elevation: \url{https://earthexplorer.usgs.gov}
\end{itemize}

No special access, permits, or funding required. Total
computational cost: approximately 10~CPU-hours in Google
Earth Engine free tier. The geometry either resonates
or it does not.

\begin{thebibliography}{9}
\bibitem{BCT_L116} M.~R.~Cabri\'{e},
BCT Letter 116: The Lost Archive,
\href{https://zenodo.org/search?q=cabrie}
{zenodo.org/search?q=cabrie} (2026).

\bibitem{BCT_L102} M.~R.~Cabri\'{e},
BCT Letter 102: BCT SAR Victorian Goldfields,
\href{https://zenodo.org/search?q=cabrie}
{zenodo.org/search?q=cabrie} (2026).

\bibitem{Mataha2008} L.~De Cordier \textit{et al.},
Mataha Expedition Report, NRIAG/Ghent (2008).

\bibitem{BEDMAP3} BAS BEDMAP3 Consortium,
\url{https://www.bas.ac.uk/project/bedmap} (2023).
\end{thebibliography}

\end{document}
